\section{Poissonization: randomize a total to reveal independence}
\label{sec:poissonization}

With exactly $m\ge1$ draws from $n\ge2$ positive-probability coupon
types, the counts add to $m$ and are dependent. Poissonization replaces
that fixed total by an independent
Poisson total. The resulting independence is exact for the new experiment;
returning to a fixed total requires conditioning or a separate comparison.
The argument also applies to balls in boxes, requests sent to servers,
and other independent categorical choices.

\subsection{Split a Poisson count}

Let $K\sim\Pois(t)$, $t>0$, independently of an i.i.d. sequence of
type labels with probabilities $p_1,\ldots,p_n>0$. Among the first $K$
labels, let $C_i$ count type $i$. For nonnegative integers $c_i$ with
sum $s$, conditioning on the total gives
\begin{align*}
\P(C_1=c_1,\ldots,C_n=c_n)
&=\frac{e^{-t}t^s}{s!}\frac{s!}{\prod_i c_i!}\prod_i p_i^{c_i}\\
&=\prod_i e^{-tp_i}\frac{(tp_i)^{c_i}}{c_i!}.
\end{align*}
Thus all the $C_i$ are independent Poisson variables with means $tp_i$.
This is the \emph{Poisson splitting} or \emph{thinning} identity.
Conversely, conditioning these independent counts on $\sum_i C_i=m$
recovers the multinomial probabilities
$m!\prod_i p_i^{c_i}/\prod_i c_i!$. For $n\ge2$ and $m\ge1$,
conditioning reintroduces dependence. With one type or zero draws,
the fixed-total count vector is constant instead.

Let $M_K$ count the missing types. The missing indicators are now
independent Bernoulli variables with probabilities $e^{-tp_i}$, so
\[
\E[M_K]=\sum_i e^{-tp_i},\qquad
\Var(M_K)=\sum_i e^{-tp_i}(1-e^{-tp_i}),\qquad
\boxed{\P(T\le K)=\prod_i(1-e^{-tp_i}).}
\]
The missing count is binomial for uniform probabilities,
$M_K\sim\Bin(n,e^{-t/n})$; with unequal probabilities it is generally
a sum of Bernoulli variables with different parameters.

\begin{example}[A load limit for every box]
Send a Poisson$(t)$ number of independent requests to boxes with routing
probabilities $p_i$. For integer $b\ge0$, independence gives
\[
\P(\text{every box has at most }b\text{ requests})
=\prod_i\left[e^{-tp_i}\sum_{j=0}^b\frac{(tp_i)^j}{j!}\right].
\]
The expected occupied-box count is $\sum_i(1-e^{-tp_i})$.
With a fixed total $m$ its expectation instead is
$\sum_i[1-(1-p_i)^m]$, by linearity; the corresponding occupancy
indicators generally are dependent.
\end{example}

\subsection{Exponential clocks for a collection}

Give the original label sequence independent time gaps
$E_1,E_2,\ldots\sim\text{Exponential}(1)$, also independent of all labels.
The arrival times $S_m=\sum_{j=1}^mE_j$ construct a rate-one Poisson
stream, by Section~\ref{sec:gamma-waits}. It also has independent counts
in disjoint intervals. Indeed, conditional on the completed arrival
history up to a fixed time, memorylessness leaves a fresh rate-one
exponential residual gap, and subsequent gaps are still independent.
Repeating this argument at successive fixed interval boundaries proves
the independent-increment assertion.

Mark each arrival with its coupon label. Splitting in each interval
makes the type-$i$ stream have rate $p_i$. To check that their first
arrival times $Y_i$ really are independent, fix thresholds $t_i\ge0$
and subdivide $[0,\max_i t_i]$ at those thresholds. In each subinterval,
the event $Y_i>t_i$ forbids labels whose threshold has not yet passed.
The probability of zero forbidden arrivals is the exponential of minus
the subinterval length times the sum of those labels' probabilities.
Independence between intervals therefore gives
\[
\P(Y_i>t_i\text{ for all }i)=\exp\!\left(-\sum_i p_i t_i\right)
=\prod_i e^{-p_it_i}.
\]
This joint survival factorization proves that
$Y_1,\ldots,Y_n$ are independent exponentials of rates $p_1,\ldots,p_n$.
The time $\tau$ when the collection is complete is their maximum, hence
\[
\P(\tau\le t)=\prod_i(1-e^{-tp_i}),\qquad t\ge0.
\]

There is an exact bridge back to the \emph{number} of draws. Since
$T$ depends only on labels and all the gaps are independent of those labels,
\[
\tau=\sum_{j\ge1}E_j\ind_{\{T\ge j\}},\qquad
\E[\tau]=\sum_{j\ge1}\E[E_j]\P(T\ge j)=\E[T].
\]
The interchange follows by truncating the nonnegative sum and taking an
increasing limit. This proves the equality of means directly; no optional
stopping theorem is being assumed. With arrival rate $\lambda$, the
same argument gives $\E[\tau]=\E[T]/\lambda$.

For a nonnegative variable, integrating its tail gives its expectation:
write its value as $\int_0^\infty\ind_{\{\tau>t\}}\,dt$ and average
nonnegative areas. Thus
\begin{align*}
\E[T]=\E[\tau]
&=\int_0^\infty\left[1-\prod_i(1-e^{-tp_i})\right]\,dt\\
&=\sum_{\varnothing\ne A}(-1)^{|A|+1}\int_0^\infty e^{-tp_A}\,dt
=\sum_{\varnothing\ne A}\frac{(-1)^{|A|+1}}{p_A}.
\end{align*}
This recovers the unequal-coupon formula. For uniform probabilities,
the substitution $u=e^{-t/n}$ instead gives
\[
\E[T_n]=n\int_0^1\frac{1-(1-u)^n}{u}\,du
=n\sum_{j=0}^{n-1}\int_0^1(1-u)^j\,du=nH_n.
\]

\begin{warningbox}
Equal means do not make the clock time and the draw count identically
distributed. Conditional on $T=r$, the clock time is Gamma$(r,1)$,
with second moment $r(r+1)$. Therefore
$\Var(\tau)=\Var(T)+\E[T]$. For one coupon, $T=1$ deterministically,
whereas $\tau$ is exponential. Treating their variances as equal would
lose the randomness added by the clock.
\end{warningbox}

\subsection{Why uniform probabilities minimize the mean}

For fixed $t>0$, the function $g(x)=\log(1-e^{-tx})$ on $x>0$ has
\[
g''(x)=-\frac{t^2e^{tx}}{(e^{tx}-1)^2}<0.
\]
It is strictly concave. The concavity inequality for a finite average
gives $n^{-1}\sum_i g(p_i)\le g(1/n)$, so
\[
\prod_i(1-e^{-tp_i})\le(1-e^{-t/n})^n.
\]
For completeness, that finite-average inequality follows by repeated
application of the two-point chord inequality for a concave function;
negative second derivative proves the chord inequality by elementary
calculus. If $n\ge2$, equality holds exactly when all $p_i=1/n$.
Comparing the integrated survival functions and using the equality of
clock and draw means proves
\[
\boxed{\E[T]\ge nH_n,}
\]
strictly for a nonuniform vector. The comparison just proved is a
stochastic comparison of \emph{clock times}, together with a comparison
of draw-count means. It does not, by itself, establish a comparison of
fixed-draw completion probabilities.

\subsection{Return to a fixed total: de-Poissonization}
\label{sec:depoissonization}

Write $F(m)=\P(T\le m)$ and $\widetilde F(t)=\P(T\le K_t)$ for an
independent $K_t\sim\Pois(t)$. Conditioning gives the exact transform
\[
\widetilde F(t)=e^{-t}\sum_{m\ge0}F(m)\frac{t^m}{m!}.
\]
For two equally likely types, $F(2)=1/2$, whereas
$\widetilde F(2)=(1-e^{-1})^2$. Replacing $K_t$ by its mean is not an
exact operation.

Monotonicity provides a rigorous comparison. For an integer $m\ge1$
and $0<t_-<m<t_+$, couple both experiments to the same label sequence.
The pointwise inequalities
$F(K_{t_-})\le F(m)+\ind_{\{K_{t_-}>m\}}$ and
$F(m)\le F(K_{t_+})+\ind_{\{K_{t_+}<m\}}$ give
\[
\boxed{\widetilde F(t_-)-\P(K_{t_-}>m)
\le F(m)\le
\widetilde F(t_+)+\P(K_{t_+}<m).}
\]
Intersect these bounds with $[0,1]$. Exact Poisson tails can be used,
or Chebyshev bounds them by $t_-/(m-t_-)^2$ and
$t_+/(t_+-m)^2$, respectively. The comparison states the error rather
than silently equating two sampling models.

\begin{example}[The completion window, as optional asymptotics]
For uniform coupons, at clock time $t=n(\log n+c)$, fixed real $c$,
\[
\widetilde F(t)=\left(1-\frac{e^{-c}}n\right)^n
\longrightarrow \exp(-e^{-c}).
\]
This holds for all sufficiently large $n$, when $t>0$.
It also describes the fixed-draw experiment. Set
$m=\lfloor n\log n+cn\rfloor$ and $t_\pm=m\pm n^{3/4}$.
Both Poisson-tail errors above tend to zero, since they are at most
$(m+n^{3/4})/n^{3/2}$, of order $\log n/\sqrt n$.
Also $t_\pm/n-\log n\to c$, so both product bounds tend to the same
$\exp(-e^{-c})$. Squeezing gives
$\P(T_n\le\lfloor n\log n+cn\rfloor)\to\exp(-e^{-c})$.
Poissonization made the coordinates independent; de-Poissonization
justified transferring the limit back to a fixed number of draws.
\end{example}
